\documentclass[submit]{ipsj}
% [submit] をつけないとフォントの問題が起きる可能性がある（jsample-gpw.tex より）

% 受付日などを消す
\makeatletter
\def\@uketsuke{}
\def\@euketsuke{}
\def\SHUBETUname{}
\makeatother

% ページ番号を消す
\pagestyle{empty}

\usepackage[dvipdfmx]{graphicx}
\usepackage[dvipdfmx]{xcolor}
\usepackage{booktabs}
\usepackage{url}
\usepackage{comment}

% 図が多いので，1 ページに置ける図の数と割合の上限を緩める
\renewcommand{\topfraction}{0.95}
\renewcommand{\dbltopfraction}{0.95}
\renewcommand{\bottomfraction}{0.6}
\renewcommand{\textfraction}{0.05}
\renewcommand{\floatpagefraction}{0.7}
\renewcommand{\dblfloatpagefraction}{0.7}
\setcounter{topnumber}{4}
\setcounter{dbltopnumber}{3}
\setcounter{bottomnumber}{2}
\setcounter{totalnumber}{6}

\newcommand{\todo}[1]{\textcolor{red}{\textbf{[未定：#1]}}}
% GUI で保存した図．ファイルがなければ枠だけ出す
\newcommand{\guifig}[3]{\IfFileExists{#2}{\includegraphics[width=#1]{#2}}%
  {\fbox{\parbox{\dimexpr#1-2\fboxsep-2\fboxrule}{\centering\small\vspace{12mm}\texttt{\detokenize{#2}}\\#3\vspace{12mm}}}}}

\title{ResTNet による将棋 AI の解釈:\\
Transformer の Attention 機構による着目領域の可視化と妥当性の検証}
\etitle{Interpreting a Shogi AI Built on ResTNet:\\
Visualizing and Validating Attention over the Board}

\affiliate{KUT}{高知工科大学情報学群\\
School of Informatics, Kochi University of Technology}

\author{横山 凜}{Rin Yokoyama}{KUT}[305151a@gs.kochi-tech.ac.jp]
\author{魏 廷翰}{Ting-Han Wei}{KUT}[tinghan.wei@kochi-tech.ac.jp]

\begin{document}

\begin{abstract}
ResTNet は Residual ブロックと Transformer ブロックを組み合わせた AlphaZero のネットワークであり，
囲碁と Hex で Attention Map による解釈が示されている．
本研究は ResTNet を動物将棋と 9×9 の将棋に適用し，Attention Map が判断の説明に使えるかを，入力が出力にどれだけ効いたかを示す Integrated Gradients（IG）で補強しながら確かめた．
完全解析された動物将棋では，Transformer ブロックだけの 6T の 1 つのヘッドが，無作為に選んだ局面で，相手のライオンを取れるときの 90\% で相手のライオンに，
負けでない局面で最善手がライオンを動かす手のときの 88\% でその行き先に，最大の重みを向け，その上位のマスは IG の上位のマスとも偶然より多く重なった．
一方，Residual ブロックを含む構成では重みがほぼ均等であった．
9×9 の将棋のプロの対局の局面では，R3(RRT) の重みはほぼ均等で，図ごとに正規化した Attention Map の濃淡は実際に指された手と対応せず，手の確率を押し上げた駒（持ち駒の桂など）は IG で読み取れた．
持ち駒は，1 枚を相手に渡したときの Attention の変化と IG で可視化した．
\end{abstract}

\begin{jkeyword}
動物将棋，ResTNet，Attention Map，Integrated Gradients，完全解析，説明可能性
\end{jkeyword}

\begin{eabstract}
ResTNet combines residual and Transformer blocks in an AlphaZero network, and its attention maps
have been used to interpret play in Go and Hex. We apply ResTNet to Dobutsu Shogi and 9×9 shogi and examine whether its attention maps explain its decisions,
checking them with Integrated Gradients (IG), which shows how much each input contributes to an output.
In the solved game of Dobutsu Shogi, one head of the all-Transformer network 6T puts its largest weight on the opponent lion in 90\% of random positions where it can be captured,
and on the destination of the best move in 88\% of the positions not lost where the best move is a lion move, and the squares it weighs most overlap with those of the largest IG more often than by chance.
In the configurations containing residual blocks, the weights are nearly uniform.
In positions of a professional game of 9×9 shogi, the weights of R3(RRT) are nearly uniform, the shading of the per-figure normalized attention maps does not match the moves played,
and the pieces that raise the probability of the move (such as a knight in hand) are found by IG.
For the pieces in hand, we show the change of attention when one piece is given to the opponent and their IG.
\end{eabstract}

\begin{ekeyword}
Dobutsu Shogi, ResTNet, attention map, Integrated Gradients, solved game, explainability
\end{ekeyword}

\maketitle
% ページ番号を消す
\thispagestyle{empty}

\section{はじめに}
AlphaZero~\cite{alphazero}は，自己対局による強化学習のみで囲碁・チェス・将棋において超人的な棋力に達した．
一方で，その判断の理由を人間が読み取ることは難しい．
判断の根拠を説明する説明可能 AI（XAI）は，医療診断や自動運転のように判断の説明が求められる分野で必要とされている．
ゲームは XAI の手法を確かめる場に適している．定跡や手筋のように言葉になった知識と照らし合わせることができ，同じ局面を何度でも再現できるからである．
囲碁や将棋のネットワークについても，その判断の根拠を可視化する研究が行われている~\cite{mandai2018,ota2024}．

ResTNet~\cite{restnet}はResidual ブロック（R）と Transformer ブロック（T）を組み合わせて並べたネットワークである．同論文は，9×9 囲碁において 6 ブロックの 11 構成を比較し，KataGo に対する勝率が最も高いのは RRTRRT であったと報告している．
また 19×19 囲碁では，10 ブロックの 10R と R3(RRT)（RRRTRRTRRT）を比較し，KataGo に対する勝率は R3(RRT) のほうが高かったと報告している．
さらに同論文は，T ブロックの attention の値から Attention Map を作成し，囲碁の死活や征（シチョウ），Hex の接続といった概念を捉えていることを示した．
ただし，同論文の解釈は個々の局面の図によるもので，重みがそれらの概念に向くかを多くの局面で数えてはいない．
また，同論文で Attention Map による解釈が示されたのは囲碁と Hex だけであり，我々の知る限り，ResTNet を将棋系のゲームに適用した例はない．
将棋系のゲームは，持ち駒・成り・駒ごとに異なる動きを持ち，行動空間も大きい（9×9 の将棋で 11,259）．
盤上の石の配置だけを扱う囲碁や Hex と比べて，盤に写らない持ち駒をどう可視化するかという問題も加わる．

図ごとに [0, 1] に正規化した Attention Map は，重みがほぼ均等でも濃淡を示す（\ref{sec:shogi}節）．
そのため，図の濃淡が判断の説明になっているかは，正解と照らし合わせなければわからない．
本研究の目的は，Attention Map を判断の説明として使えるかを，正解のわかるゲームで確かめることである．
ここでは，Attention Map が説明になっているとは，局面の正解に関わるマス（取れる相手のライオンや，最善手の行き先など）に，選んだ例だけでなく無作為な局面でも重みが向くことを指す．
将棋系のゲームは，この確かめ方を将棋の知識の解釈に広げる前の土台として用いる．

動物将棋は完全解析されており，初期局面は 78 手で後手勝ちである~\cite{tanaka2009}．
したがって，全局面の最善手と勝敗を正解として参照でき，Attention Map が正解のマスを見ているかを確かめられる．
一方，9×9 の将棋には正解がない．9×9 の将棋では，Attention Map の可視化に利用した棋譜について，実際に指された手を仮の正解として，Attention Map が将棋としてそれらしいかを見るにとどめる．

本研究では，ResTNet を将棋系のゲームに適用し，次を行った．
    \begin{itemize}
        \item 動物将棋の 11 構成を自己対戦で学習し，R ブロックだけの構成と T ブロックを含む構成で最善手を選ぶ割合を比べた．
        \item 動物将棋では，Attention Map の重みが向くマスが完全解析の正解と合うかを，ヘッドを選ぶ局面とは別の，完全解析から無作為に選んだ局面で確かめた．
        \item Attention Map が重みを向けたマスが出力にも効いているかを，Integrated Gradients（IG）で確かめた．
        \item 9×9 の将棋では，Attention Map を図ごとに [0, 1] に正規化する前の重みの偏りも示し，正規化が均等に近い重みにも濃淡を作ることを示した．
        \item Attention Map に写らない持ち駒について，持ち駒を 1 枚相手に渡したときの Attention の変化と，持ち駒の IG を可視化した．
    \end{itemize}

\section{背景と関連研究}
\subsection{AlphaZero と Gumbel AlphaZero}
AlphaZero~\cite{alphazero} は，深層ニューラルネットワークとモンテカルロ木探索（MCTS）を組み合わせ，
人間の知識を用いずに自己対戦だけで学習する手法である．
ネットワークは残差ブロックを重ねた本体と，policy ヘッド，value ヘッドからなる．
policy ヘッドは各手の確率の分布を出し，MCTS の根での各手の訪問回数の分布に近づくように学習する（Gumbel AlphaZero では，探索で改善した方策に近づくように学習する）．
value ヘッドは局面の勝ちやすさを出し，対局の結果に近づくように学習する．
MCTS は，policy を各手の事前確率として，value を局面の評価として用いる．
学習では，MCTS による自己対戦で棋譜を作ることと，その棋譜でネットワークを最適化することを繰り返す．
Gumbel AlphaZero~\cite{gumbel} は，少ないシミュレーション数でも方策の改善が保証されるように AlphaZero を改良した手法である．

\subsection{ResTNet}
\label{sec:restnet}
ResTNet~\cite{restnet} は，R ブロックと T ブロックを並べたネットワークである．
R ブロックは AlphaZero と同じく，2 層の畳み込みからなる残差ブロックである．
T ブロックは，盤の各マスを 1 トークンとし，相対位置の表現を加えた Transformer のブロックである．
同論文は，9×9 囲碁で R と T の並べ方が異なる 6 ブロックの 11 構成を比べ，19×19 囲碁と 19×19 Hex で 10 ブロックの 10R と R3(RRT) を比べた．

T ブロックでは，各マスの特徴から query・key・value の 3 つのベクトルを作る．
あるマス（以下，起点）の query と各マスの key の内積に相対位置のバイアスを足して softmax をとったものが attention の重みで，0〜1 の値をとり，起点ごとに全マスの合計が 1 になる．
起点の新しい特徴は，この重みで各マスの value ベクトルを足し合わせたもの（重みの合計が 1 なので加重平均）であり，これを入力に足して（残差）次の処理に渡す．
同論文は，起点が盤の各マスに向ける attention の重みを盤の形に並べた図を Attention Map とし，重みを「ほかのマスの相対的な重要さ」としている．
同論文は 19×19 囲碁の R3(RRT) で，生き残る石（死活），どちらの地にもなっていない中央，シチョウで逃げる石とつながる石に，
19×19 Hex で石の接続に，重みが集まる図を示した．値は [0, 1] に正規化している．

\subsection{ゲーム AI の判断の説明}
万代と金子~\cite{mandai2018} は，画像分類で用いられる Saliency Map と SmoothGrad を囲碁の policy と value に適用し，予測に大きく影響する座標を示した．
大田と中沢~\cite{ota2024} は，将棋 AI の dlshogi が policy と value を出すときに重視する駒の利きを可視化した．
Integrated Gradients（IG）~\cite{ig} は，基準の入力から実際の入力まで入力を少しずつ変えたときの出力の勾配を積分し，各入力が出力にどれだけ効いたかを求める手法である．各入力の値の合計は，出力と基準の入力での出力の差に等しい．

\subsection{動物将棋}
動物将棋は 3×4 の盤で，ライオン・きりん・ぞう・ひよこの駒を用いる将棋である~\cite{tanaka2009}．
取った駒は持ち駒として打てる．相手のライオンを取るか，自分のライオンが相手の陣地に入り，直後に取られなければ（トライ）勝ちである．
完全解析されており~\cite{tanaka2009}，すべての局面の勝敗と最善手がわかるため，本研究では Attention Map を正解と照らし合わせるために用いた．

\subsection{9×9 の将棋}
9×9 の盤で行う通常の将棋である．
行動空間は AlphaZero と同じ 11,259（9×9×139）である~\cite{alphazero}．
本研究では 512 手に達した対局は引き分けとし，入玉宣言（27 点法）を有効とした．
本文の指し手は，▲が先手，△が後手の手を表す（例：▲5五銀は先手が銀を 5五 に動かす手，△5五銀打は後手が銀を 5五 に打つ手，動物将棋の ▲キ a3 打は先手がきりんを a3 に打つ手）．

\section{手法}
\subsection{ネットワークと学習}
ResTNet の実装~\cite{restnet} を用い，動物将棋と 9×9 の将棋で同じ形のブロックを使う．
どちらのゲームでも，入力を 3×3 の畳み込みで 256 枚の平面に変換し，以降の R ブロックと T ブロックはすべてこの 256 枚の上で計算する．
R ブロックの畳み込みは 3×3，T ブロックは 4 ヘッドである．policy ヘッドは AlphaZero と同じ畳み込み型，value ヘッドは Transformer 型とした．
2 つのゲームで異なるのは，盤の大きさ（3×4 と 9×9），入力の枚数（\ref{sec:input}節），policy が出す手の数（132 と 11,259），ブロックの数と並べ方である．
動物将棋は 6 ブロックで，ResTNet 論文の 9×9 囲碁と同じ 11 構成
（6R, 6T, 5R1T, 4R2T, 3R3T, 2R4T, 1R5T, TRRRRT, RTRRRT, RRTRRT, RRRTRT）を学習した．
9×9 の将棋は，ResTNet 論文が 19×19 囲碁で用いた 10 ブロックの R3(RRT)（RRRTRRTRRT）を学習した．

学習には MiniZero~\cite{minizero} 上の Gumbel AlphaZero~\cite{gumbel} を用いた（表~\ref{tab:train}）．
自己対戦で指す手は，根の子のうち，最も訪問回数の多い子との価値の差が 0.1 以内のものから，訪問回数に比例して抽選した（MiniZero の初期設定）．
9×9 の将棋のシミュレーション数，Gumbel のサンプル数，バッチサイズ，学習率とその半減，自己対戦の局数（100 万局），最適化のステップ数（100,000）は，ResTNet 論文の 9×9 囲碁と同じである．

\begin{table}[tb]
  \centering\footnotesize
  \caption{学習の設定}
  \label{tab:train}
  \begin{tabular}{lll}
    \toprule
    & 動物将棋 & 9×9 の将棋 \\
    \midrule
    シミュレーション数 & 64 & 64 \\
    Gumbel のサンプル数 & 16 & 16 \\
    1 iteration の自己対戦 & 2,000 局 & 2,000 局 \\
    1 iteration の最適化 & 200 ステップ & 200 ステップ \\
    iteration 数 & 200 & 500 \\
    バッチサイズ & 1,024 & 1,024 \\
    最適化手法 & SGD & SGD \\
    学習率（初め） & 0.02 & 0.02 \\
    学習率を半分にするステップ & 28,000，36,000 & 70,000，90,000 \\
    学習率（最後） & 0.005 & 0.005 \\
    リプレイバッファ & 20 iteration & 15 iteration \\
    投了 & 無効 & 閾値 $-0.9$ \\
    \bottomrule
  \end{tabular}
\end{table}

\subsection{入力}
\label{sec:input}
入力は AlphaZero~\cite{alphazero} と同じく，盤と同じ大きさの平面を重ねたものである．
盤上の駒は，手番側と相手のそれぞれについて駒の種類ごとに 1 枚の平面で表し，その駒があるマスを 1 とする．
持ち駒は駒の種類ごとに 1 枚の平面で表し，盤の全マスにその枚数を入れる．
9×9 の将棋は AlphaZero と同じく，同じ局面が現れた回数の平面（3 枚）を加えた直近 8 局面の平面を重ね，手番と手数の平面を加える（362 枚）．
動物将棋は現在の局面だけを入力する（20 枚）．

\subsection{Attention Map}
\label{sec:attn}
T ブロックの attention の重み（softmax の後）を用いる．
起点は任意のマスに指定できる．本稿では，動物将棋はモデルが選んだ手（図~\ref{fig:concepts}(d) を除く），9×9 の将棋は実際に指された手で動かす駒のマスを起点とし，駒打ちでは打つ先のマスを起点とした．
Attention Map は T ブロックにしかないため，ResTNet の論文と同じく，T ブロックだけを入力側から数えて「第 $n$ T ブロック」と呼ぶ．
各 T ブロックには 4 つのヘッドがあり，「第 $h$ ヘッド」と呼ぶ．
6T はすべてのブロックが T ブロックなので，第 $n$ ブロックがそのまま第 $n$ T ブロックである．
RRTRRT では第 3・6 ブロックが第 1・2 T ブロック，R3(RRT) では第 4・7・10 ブロックが第 1・2・3 T ブロックである．

重みの偏りは，重みにマスの数 $N$ を掛けた値で表す（動物将棋は $N = 12$，9×9 の将棋は $N = 81$）．
この値は，起点が全マスに同じ重みを向けたとき全マスで 1，1 マスだけに向けたときそのマスで $N$ になる．
動物将棋の図では，この値を全部の図で同じ目盛りで描く．
9×9 の将棋の重みはほぼ均等であった（\ref{sec:shogi}節）．そのため，ResTNet と同じく図ごとに最小の重みを 0，最大の重みを 1 に引き伸ばして描き，図の上に重み × 81 の最大を書く．
この表示では，均等に近い重みにも濃淡が現れるため，図の濃淡は最大の値と合わせて読む必要がある．

重みがほぼ均等であることは，T ブロックが何もしていないことを意味しない．
\ref{sec:restnet}節のとおり，起点の新しい特徴は各マスの value ベクトルの重み付きの平均であり，重みが均等なら盤全体の平均になる．

図には，Attention Map と比べるために，モデルが選んだ手を示す．この手は探索を用いず，policy が最大の合法手とする．
Attention Map は 1 回の推論の中の値であり，同じ推論の出力である policy と比べるためである．
対局で指す手は MCTS で選ぶが，MCTS の手はほかの多くの局面の推論の結果も反映するため，1 つの局面の Attention Map とは直接は対応しない．

\subsection{Integrated Gradients}
\label{sec:ig}
Attention Map は内部の重みであり，モデルの出力（手の確率や value）がどの入力で決まったかは示さない．
そこで，Attention Map が重みを向けたマスが出力にも効いているかを，IG~\cite{ig} で確かめる．
IG は，全部 0 の入力（盤上の駒も持ち駒もなく，手番などの平面も 0）を基準とし，基準から今の局面まで入力を少しずつ増やしたとき，それぞれの入力の分として出力がどれだけ動いたかを表す（300 分割の和で近似）．
値の単位は出力と同じで，全部の入力の値を足すと「今の局面の出力 − 基準での出力」になる．
正の値はその駒が出力を押し上げたこと，負の値は押し下げたことを表す．ある駒 1 枚を取り除いたときの出力の変化ではない．
出力には，value（手番側から見た勝ちやすさ）と，モデルが選んだ手（9×9 の将棋は実際に指された手）の確率を用いる．
図では，盤上の駒の値をマスごとに足して盤の形に描き（9×9 の将棋は直近 8 局面分を足す），押し上げた駒を赤，押し下げた駒を青で示す（目盛りは図ごと）．
持ち駒は盤に場所がないので，駒の種類ごとの棒で示す．

\subsection{持ち駒の可視化}
持ち駒は盤の全マスに同じ値を入れた平面で入力されるため，Attention Map には持ち駒の場所がない．
そこで，最善手（9×9 の将棋は実際に指された手）が駒打ちの局面で，次の 3 つを見る．

\textbf{持ち駒を 1 枚相手に渡したときの Attention の変化．}
打つ先を起点とし，打つ駒を持っているときの重み × マスの数から，その駒を 1 枚相手の持ち駒に移したとき（駒の総数は変わらない）の重み × マスの数を引いた差を盤の形に描く．
重みは入力全体から計算されるため，持ち駒の平面が変わると重みも変わる．
赤いマスはその持ち駒を持っているときに多く重みを向けるマス，青いマスは少なく向けるマス，白は変わらないマスで，目盛りは全部の図で同じ（$\pm 7$）である．

\textbf{持ち駒の IG．}
打つ手の確率を，持ち駒がどれだけ押し上げたかを IG の棒で見る．

\textbf{打つ確率．}
打つ駒の種類ごとに，各マスに打つ手の policy の確率を盤の形に並べ，最善手（9×9 の将棋は実際に指された手）の打つ先に青の○，モデルが選んだ手の打つ先に黒い点線の○を付ける．

\subsection{ヘッドと局面の選び方}
\label{sec:select}
\textbf{動物将棋のヘッド．}
ヘッドを選ぶ局面と，選んだヘッドを確かめる局面を分けた．
選ぶ局面には，完全解析の最善手順の 61 局面（判定した手 71 個）を用いた．
各局面で最善手が動かす駒のマスを起点とし，行き先のマスへの重み × 12 を，T ブロックを持つ 10 構成の全 116 ヘッドについて平均した（図~\ref{fig:select}）．
均等（1）を明確に上回ったのは 6T の第 5 T ブロックの第 4 ヘッドだけで（1.99），2 番目は同じ 6T の第 2 T ブロックの第 1 ヘッドの 1.12 であった．
6T 以外の 9 構成の 92 ヘッドは，すべて 0.994〜1.006 であった．
このヘッドを例の図に用い，選ぶのに使わなかった無作為な局面で確かめた（\ref{sec:dobutsu}節）．
6T の第 4 T ブロックの第 1 ヘッドは，例の局面を見る中で重みが大きく偏っていたヘッドで，図~\ref{fig:concepts}(d) の見方（相手のライオンから自分のライオンに重みを向ける）は無作為な局面では確かめていない．
なお，ResTNet の論文も図ごとにどのブロック・ヘッドかを示しているが，選び方は書かれていない．

\begin{figure}[tb]
  \centering
  \guifig{\columnwidth}{figs/gui_d_validity.png}{ヘッドの選び方の図}
  \caption{動物将棋：ヘッドの選び方．最善手順の 61 局面で，最善手が動かす駒のマスから行き先のマスへの重み × 12 を平均した値（点 1 つが 1 ヘッド，T ブロックを持つ 10 構成の全 116 ヘッド）．点線は均等（1）．}
  \label{fig:select}
\end{figure}

\textbf{動物将棋の局面．}
完全解析から条件ごとに選んだ候補の中から，説明しやすい局面を選んだ．
候補は，図~\ref{fig:dexample} が 6T が勝ちを失った局面から無作為に選んだもの（そのうち 6T の重みが取れる相手のライオンに集まっていた局面），図~\ref{fig:concepts}(a) が 6T が最善手の詰みを外した局面から無作為に選んだもの，
図~\ref{fig:concepts}(b)〜(d) がそのヘッドの重み × 12 の最大（5 以上）が正解のマスに向いた局面から無作為に選んだもの，
図~\ref{fig:hand} が持ち駒を 1 枚相手に渡したときの重みの変化が最も大きい局面である．後の 2 つは，示す性質そのもので選んだ局面である．
したがって，これらの図は例示であり，選ばなかった局面でも同じことが成り立つかは示さない．
それは\ref{sec:dobutsu}節で，完全解析から無作為に選んだ局面で，ヘッドの重みが最大となるマスが正解のマスと一致する割合を数えて確かめる．

\textbf{9×9 の将棋の局面．}
正解がないため，プロの対局の棋譜を用い，実際に指された手を仮の正解とした．
藤井聡太五段と羽生善治竜王の対局（第 11 回朝日杯将棋オープン戦 準決勝，2018 年）\footnote{棋譜は \url{https://sakidori-ch.com/fujiisouta-20180217}（出典：連盟モバイル）による．}から，
R3(RRT) が選ぶ手が実際に指された手と一致した 80 手目までの局面（次の 81 手目に ▲5五銀）と 103 手目までの局面（次に △5五銀打），
および最後の手の直前の 118 手目までの局面（次の 119 手目に ▲4七桂打．王手で，直後に後手が投了）の 3 局面を用いた．
図には 118 手目までの局面を示す．

\section{結果と考察}
\subsection{モデルの強さ}
\label{sec:reliability}
可視化するモデルがどの程度学習しているかを確かめる．

\textbf{動物将棋．}
モデルが選んだ手（探索なしで policy が最大の合法手）が，最善手（勝ちなら最短で勝つ手，負けなら最も長く粘る手，引き分けなら引き分けを保つ手）と同じ評価である局面の割合を「最善手を選んだ割合」とし，
完全解析の全 246,803,167 局面から，相手のライオンがトライ済みの負け局面 7,018,985（どの手も同じ評価）を除く 239,784,182 局面で数えた．
11 構成とも 0.82〜0.87 で，6T が最も高い（先手番として入力したとき 0.865，後手番として入力したとき 0.866）．
6T を先手番として入力したときの内訳は，相手のライオンを取れる局面で 0.973，それ以外の勝ちの局面（最短で勝つ手）で 0.672，引き分けの局面で 0.734，負けの局面で 0.769 であった．
すぐに勝てる局面ではほぼ必ず勝つ手を選ぶ一方，勝ちまで手数のかかる局面では，最短で勝つ手を選べるのは 3 分の 2 ほどである．
R ブロックだけの 6R は 0.843（先手番として入力したとき）で，11 構成の差は最大 0.041 であった．
動物将棋では，T ブロックを加えても最善手を選ぶ割合は大きくは変わらなかった．
これは探索なしの policy の 1 手で数えた値であり，探索を用いた対局の強さとは別の量である．

\textbf{9×9 の将棋．}
正解がないため，対局で強さを示す．
ResTNet~\cite{restnet} が KataGo の学習途中のモデルを相手にしたのにならい，AobaZero の学習初期の重み（w50, w60）と対局させた．
AobaZero は AlphaZero の将棋の実験の追試を目的としたプログラムで~\cite{aoba2020,aoba2021}，学習途中の重みを公開している\footnote{\url{http://www.yss-aya.com/aobazero/}}．
AobaZero（v18）は 1 手 800 playout，1 スレッド，乱数なしとし，こちらは 1 手 64 シミュレーション（学習と同じ），ノイズなしで，sequential halving で最後に残った候補を指す．
開始局面は出村と金子~\cite{demura2023} と同じく，人間の棋譜集の 9 戦型の 20〜40 手目から選び，各局面で先後を入れ替えて 2 局ずつ指した（各重みに 400 局）．
R3(RRT) の勝率（勝ち ÷ 局数）は，w50 に対して 60.8\%（243 勝 9 分 148 敗），w60 に対して 31.5\%（126 勝 3 分 271 敗）であった．

\subsection{動物将棋：Attention Map と IG}
\label{sec:dobutsu}
\textbf{Attention Map．}
モデルが誤った局面で Attention Map がどこに重みを向けるかを見るため，6T が勝ちを失った局面を選んだ（図~\ref{fig:dexample}）．
後手は a2a3 で相手のライオンを取れるが（青い矢印），6T は自分のライオンを c2b1 と動かした（黒い点線の矢印）．
6T の第 5 T ブロックの第 4 ヘッドは，起点 c2 から相手のライオンのいる a3 に重みを集めた（6.69）．
同じ T ブロックのほかの 3 ヘッドの最大は 1.05〜1.11 で，同じく c2b1 を選んだ RRTRRT の第 2 T ブロックの 4 ヘッドは，どれも最大 1.01 とほぼ均等であった．
重みは取れる相手のライオンに集まっていたが，policy は別の手を選んだ．

\textbf{IG による確かめ．}
同じ局面の IG（図~\ref{fig:dig}）で，a3 のライオンが出力にも効いているかを見る．
6T の value（後手から見た勝ちやすさ）は 0.99 である．
IG でこの value を最も押し下げたのは a3 の先手のライオン（$-1.45$）で，押し上げたのは後手のライオン c2（$+1.47$），後手の持ち駒のぞう（$+0.76$），a3 のライオンを取れる後手のにわとり a2 と後手のきりん c1（ともに $+0.70$）であった．
a3 のライオンの値の絶対値は盤上の駒で 2 番目に大きく，Attention Map で重みが集まったマスは value にも大きく効いている．
一方，c2b1 の確率（0.81）を最も押し上げたのは動かす後手のライオン c2 自身（$+1.89$）で，a3 のライオンは押し下げる向き（$-0.60$，盤上の駒で 4 番目）に働いた．
IG の上では，c2b1 の確率の高さは主に c2 のライオン自身で説明され，取れる相手のライオンはこの確率を押し下げる側にあった．

\begin{figure}[tb]
  \centering
  \guifig{0.48\columnwidth}{figs/gui_d_lost_6T_b5h4.png}{6T 第5Tブロック 第4ヘッド}
  \caption{動物将棋：6T が勝ちを失った後手番の局面での，第 5 T ブロックの第 4 ヘッドの Attention Map．
    起点（黒枠）は手番側のライオン c2．各マスの値は重み × 12（盤の色が 1，赤いほど大きい．最大 12）．
    青い矢印は完全解析の最善手 a2a3，黒い点線の矢印は 6T が選んだ手 c2b1（探索なしで policy が最大の合法手）．
    駒は 1 文字（ラ：ライオン，キ：きりん，ゾ：ぞう，ヒ：ひよこ，ニ：にわとり），逆さの駒が後手．}
  \label{fig:dexample}
\end{figure}

\begin{figure*}[tb]
  \centering
  \guifig{0.8\textwidth}{figs/gui_d_ig_lost.png}{IG lost}
  \caption{動物将棋：図~\ref{fig:dexample} の局面での 6T の IG．左：value（後手から見た勝ちやすさ），中：6T が選んだ c2b1 の確率．
    赤は出力を押し上げた駒，青は押し下げた駒（目盛りは図ごと）．右：持ち駒の IG（駒の種類ごとの棒）．矢印は図~\ref{fig:dexample} と同じ．}
  \label{fig:dig}
\end{figure*}

\textbf{局面の種類ごとの例．}
図~\ref{fig:concepts} は，同じヘッドが局面の種類によって違うマスに重みを集める例と，別のヘッドが別のマスに重みを集める例である
（ResTNet の論文も，ブロックやヘッドごとに違う種類の知識を捉えうるとしている）．
(a) 詰み：第 5 T ブロックの第 4 ヘッドは，起点 b4 から詰ませるマス b2 に 5.98 を集めたが，6T は b4c4 と指した．
(b) 狙ってくる駒：起点 c1 から，自分のライオンを狙う c2 の駒に 7.24 を集め，6T はその駒を取った．
(c) 粘りの手：負けの局面で，最も長く粘る逃げ先 a2 に 7.08 を集め，6T はそこへ逃げた．
(d) 第 4 T ブロックの第 1 ヘッドは，相手のライオン（c1）を起点にすると，自分のライオン（b4）に 9.50 を集めた．
トライの局面でも，6T が取られるトライ a3b4 を指した局面では，第 5 T ブロックの第 4 ヘッドはその行き先 b4 に 6.22 を集め，安全なトライの行き先 a4 には 2.33 であった．

\begin{figure}[tb]
  \centering
  \begin{minipage}[t]{0.48\columnwidth}\centering (a)\\ \guifig{\linewidth}{figs/gui_fig3_d.png}{(a)}\end{minipage}\hfill
  \begin{minipage}[t]{0.48\columnwidth}\centering (b)\\ \guifig{\linewidth}{figs/gui_fig3_e.png}{(b)}\end{minipage}\\[2mm]
  \begin{minipage}[t]{0.48\columnwidth}\centering (c)\\ \guifig{\linewidth}{figs/gui_fig3_f.png}{(c)}\end{minipage}\hfill
  \begin{minipage}[t]{0.48\columnwidth}\centering (d)\\ \guifig{\linewidth}{figs/gui_fig3_g.png}{(d)}\end{minipage}
  \caption{6T の第 5 T ブロックの第 4 ヘッド（(a)〜(c)）と第 4 T ブロックの第 1 ヘッド（(d)）が局面の種類ごとに重みを集めるマス（表示は図~\ref{fig:dexample} と同じ）．
    (a) 1 手詰めに相当する後手番の局面：最善手 b3b2・c3b2，6T は b4c4．
    (b) ライオンを狙う駒がある先手番の局面：6T は c1c2（最善手の 1 つ）．
    (c) 負けの先手番の局面：6T は最も長く粘る a1a2．
    (d) 後手番の局面：起点は相手のライオン c1．}
  \label{fig:concepts}
\end{figure}

\textbf{無作為な局面での確かめ．}
これらは説明しやすい局面を選んだ例示である．
選んでいない局面でも成り立つかを確かめるため，完全解析から無作為に選んだ局面で，手番側のライオンを起点に，各ヘッドの重みが最大となるマスを調べた（図~\ref{fig:stats}）．
第 5 T ブロックの第 4 ヘッドでは，相手のライオンを取れる局面の 90.1\% で相手のライオン，負けでない局面で最善手がライオンの手のときの 87.7\% で最善手の行き先，
ライオン以外の駒に自分のライオンが狙われている局面の 65.2\% でその駒，負けの局面で最も長く粘る手がライオンの手のときの 51.0\% でその行き先に重みが最大となった．
4 つの問いのうち 3 つでは，6T のほかの 23 ヘッドのどれよりも高かった．負けの局面で粘る手の行き先では，ほかのヘッドの最も高いもの（52.8\%）のほうが高かった．
RRTRRT の 8 ヘッドで最も高いものは 4 つの問いで 11.6，10.4，9.2，13.1\% で，偶然（9.1，9.3，9.1，15.7\%）との差は 2.6 ポイント以内であった．
T ブロックごとに見ると，同じ局面（ライオンが盤上にある 19,434 局面）で手番側のライオンを起点とした重み × 12 の最大（4 ヘッドの最大）の中央値は，
6T の第 1〜6 T ブロックで 1.50，1.41，1.21，4.49，5.16，1.22 であった．重みの偏りは深い第 4・5 T ブロックで大きく，最後の第 6 T ブロックでは小さい．RRTRRT は 2 つの T ブロックとも 1.01 であった．

IG でも同じマスが出力に効いているかを，このうち 1,000 局面で確かめた．
value の IG の絶対値が大きい上位 3 マスと，重みの大きい上位 3 マスの重なりの平均は，
6T の第 5 T ブロックの第 4 ヘッドで 1.20，第 4 T ブロックの第 1 ヘッドで 1.31 と偶然（0.75）を上回り，RRTRRT は最も重なるヘッドでも 0.74 であった．

したがって，6T の第 5 T ブロックの第 4 ヘッドが相手のライオンや最善手の行き先に重みを向けることは，選んだ例だけでなく，無作為な局面でも成り立ち，IG の上位のマスとも偶然より多く重なる．
一方で，図~\ref{fig:dexample} や図~\ref{fig:concepts}(a) のように，正解のマスに重みを向けていても policy が別の手を選ぶ局面がある．
Attention Map が示すのは，起点の特徴を作るときにどのマスの特徴を用いたかであり，選んだ手の理由そのものではない．

\begin{figure}[tb]
  \centering
  \guifig{\columnwidth}{figs/gui_fig4_stats.png}{集計の図}
  \caption{動物将棋：手番側のライオンを起点に，最大の重みが正解のマスに向いた割合（完全解析から無作為に選んだ局面，先手番・後手番として入力したもの各 10,000 局面．横軸の下は問いごとの局面の数）．
    6T の第 5 T ブロックの第 4 ヘッド，6T のほかの 23 ヘッドで最も高いもの，RRTRRT の 8 ヘッドで最も高いもの（問いごとに選ぶ），偶然（ライオン以外の 11 マスから 1 つ選んだとき）．}
  \label{fig:stats}
\end{figure}

\subsection{9×9 の将棋：Attention Map と IG}
\label{sec:shogi}
\textbf{Attention Map．}
図~\ref{fig:sexample} は，118 手目までの局面での R3(RRT) の Attention Map である．
3 局面とも，3 つの T ブロックの 12 ヘッドのうち起点からの重みの偏りが最も大きいのは第 1 T ブロックの第 2 ヘッドで，
重みに 81 を掛けた値は 0.974〜1.030 の範囲にあり，ほぼ均等であった．
T ブロックごとに見ると，起点からの重み × 81 の最大（4 ヘッドの最大）は，3 局面とも第 1・2・3 T ブロックの順に 1.02〜1.03，1.01，1.00 で，深い T ブロックほど均等に近かった．
[0, 1] に正規化した図では，3 局面とも，最も薄い 5 マスはすべて空きマスで，濃いマスの多くは駒のあるマスであった．
118 手目までの局面では，先手陣の左下の隅（後手の銀 9九・と 8九・角 7九，先手の金 9八・香 9七・歩 8七）が濃い．
一方，3 局面とも，実際に指された手の行き先や，相手の玉のマスは，濃いマスに入っていなかった．
重みの差は小さい（最大で均等の 1.03 倍）ため，この濃淡を判断の根拠と言えるかは確かめられていない．

実際に指された手を仮の正解とすると，R3(RRT) が選んだ手は，80 手目までの局面の ▲5五銀（確率 0.88）と 103 手目までの局面の △5五銀打（0.65）では実際の手と一致した．
ただし 103 手目までの局面は王手を受ける局面で，合法手は 3 つしかない．
118 手目までの局面では，R3(RRT) は ▲5六銀打を選び，▲4七桂打には 0.09 を置いた．

重みが均等に近くなる理由は確かめていないが，考えられる理由の 1 つは R ブロックとの役割の分担である．
動物将棋でも，T ブロックだけの 6T の一部のヘッドは重みが大きく偏る一方，R ブロックを含む 9 構成のヘッドはどれもほぼ均等であった（図~\ref{fig:select}）．
R ブロックが近くのマスどうしの形を処理し，T ブロックは盤全体の平均を全マスに配る役割（\ref{sec:attn}節）になった可能性がある．

\textbf{IG による確かめ．}
同じ局面の IG（図~\ref{fig:sig}）を見る．
R3(RRT) が ▲4七桂打に置いた確率は 0.09 である．IG でこの確率を最も押し上げたのは持ち駒の桂で（$+0.06$），盤上の駒の値はどれも 0.03 以下であった．
IG の上では，この手の確率は，桂を持っていることで説明される．
value（先手から見た勝ちやすさ）では，盤上の駒では先手のと 2二（$+0.20$）が最も押し上げ，後手のと 7七（$-0.16$）が最も押し下げた．
持ち駒では，先手の銀・金が押し上げ（$+0.14$，$+0.10$），後手の銀・歩が押し下げた（$-0.22$，$-0.18$）．
ほかの 2 局面では，▲5五銀の確率を最も押し上げたのは動かす先手の銀 5六（$+0.64$）で，△5五銀打の確率では王手をかけられている後手の玉 4四（$+0.48$）であった．
IG の合計と，出力と基準の入力での出力の差は，3 局面とも 0.006 以内で一致した．

動物将棋の 6T と異なり，IG と Attention Map は同じマスを指さなかった．
value と実際に指された手の確率のそれぞれについて，IG の絶対値が大きい上位 5 マスと，Attention Map の濃い上位 5 マスの重なりは，3 局面で 0〜2 マスであった（偶然の期待値は約 0.3 マス）．
9×9 の将棋の R3(RRT) では，手の確率を押し上げた駒（動かす銀，持ち駒の桂）は，Attention Map ではなく IG で読み取れた．

\begin{figure}[tb]
  \centering
  \guifig{0.32\textwidth}{figs/gui_9x9_m119_paper.png}{119}
  \caption{9×9 の将棋：藤井－羽生戦の 118 手目までの局面（次に ▲4七桂打）での，R3(RRT) の第 1 T ブロックの第 2 ヘッドの Attention Map．
    起点（黒枠）は打つ先 4七．図ごとに [0, 1] に正規化（赤いほど重みが大きい）．
    青い矢印は実際に指された手，黒い点線の矢印は R3(RRT) が選んだ手（探索なしで policy が最大の合法手）．駒打ちは盤の外から打つ先への矢印で示し，矢印の元に打つ駒を書く．}
  \label{fig:sexample}
\end{figure}

\begin{figure*}[tb]
  \centering
  \guifig{0.8\textwidth}{figs/gui_9x9_m119_igv.png}{IG value}\\[2mm]
  \guifig{0.8\textwidth}{figs/gui_9x9_m119_igp.png}{IG policy}
  \caption{9×9 の将棋：図~\ref{fig:sexample} の局面での R3(RRT) の IG．上：value（先手から見た勝ちやすさ），下：▲4七桂打の確率．
    左は盤上の駒の IG（直近 8 局面分の和．赤は押し上げた駒，青は押し下げた駒．目盛りは図ごと），右は持ち駒の IG．矢印は図~\ref{fig:sexample} と同じ．}
  \label{fig:sig}
\end{figure*}

\subsection{持ち駒}
\textbf{動物将棋．}
図~\ref{fig:hand} は，最善手がきりんを a3 に打つ手（▲キ a3 打）で，6T も同じ手を選ぶ局面である．
a3 に打ったきりんは，隣の b3 にいる後手のライオンを取れる位置にある．
6T の第 4 T ブロックの第 1 ヘッドでは，きりんを持っているときのほうが，起点 a3 から c3 に向ける重み × 12 が 2.22 大きく，c2 に向ける値が 5.73 小さかった．
きりんの有無で，打つ先の特徴を作るときに用いるマスが変わったことになる．
ただし，その変化が ▲キ a3 打を選ぶことにどう効くかは，この図からはわからない．
そこで同じ局面の IG を見る（図~\ref{fig:hig}）．
6T は ▲キ a3 打に 0.79 の確率を置いた．IG でこの確率を最も押し上げたのは持ち駒のきりん（$+1.08$）であった．
値が確率 0.79 より大きいのは，盤上の駒には押し下げたもの（a1 の先手のきりんの $-0.49$ など）もあり，全部を足すと「0.79 − 基準での確率」になるためである．
IG の上では，この手の確率の高さは，きりんを持っていることで説明される．
value（先手から見た勝ちやすさ，0.98）では，後手のライオン b3 が押し上げ（$+0.62$），先手のぞう b2 が押し下げた（$-0.55$）．

\textbf{9×9 の将棋．}
103 手目までの局面（△5五銀打）と 118 手目までの局面（▲4七桂打）で，打つ駒を 1 枚渡しても，起点からの重み × 81 の変化は 0.001 以下であった．
9×9 の将棋の R3(RRT) では，持ち駒の有無は Attention Map にほとんど表れなかった．
一方，IG では，▲4七桂打の確率を持ち駒の桂が最も押し上げた（\ref{sec:shogi}節）．
打つ確率（図~\ref{fig:hdrop}）では，118 手目までの局面は，桂を打つ確率は 4七 が最も高い（0.09）が，駒打ち全体では銀を 5六 に打つ手（0.24）が最も高く，R3(RRT) はこの手を選んだ．銀を 6六 に打つ手も 0.17 であった．
103 手目までの局面は，銀を 5五 に打つ手に 0.65 を置いた．

\begin{figure}[tb]
  \centering
  \guifig{0.32\textwidth}{figs/gui_d_hand_diff.png}{diff}
  \caption{動物将棋：最善手が ▲キ a3 打の先手番の局面（先手はきりんを 1 枚持つ）での，6T の第 4 T ブロックの第 1 ヘッドの，起点 a3 からの重み × 12 の差
    （きりんを持っているとき − 1 枚相手に渡した後．赤は持っているときに多い，青は少ない）．駒打ちは盤の外から打つ先への矢印（元に「キ打」）で示す．}
  \label{fig:hand}
\end{figure}

\begin{figure*}[tb]
  \centering
  \guifig{0.8\textwidth}{figs/gui_d_ig_hand.png}{IG hand}
  \caption{動物将棋：図~\ref{fig:hand} の局面での 6T の IG（表示は図~\ref{fig:dig} と同じ）．左：value，中：▲キ a3 打の確率，右：持ち駒の IG．}
  \label{fig:hig}
\end{figure*}

\begin{figure*}[tb]
  \centering
  \guifig{0.56\textwidth}{figs/gui_9x9_m119_drop.png}{119 drop}
  \caption{9×9 の将棋：図~\ref{fig:sexample} の局面で R3(RRT) が桂と銀を各マスに打つ確率（青の○は実際に指された手，黒い点線の○は R3(RRT) が選んだ手の打つ先．0.01 以上のマスに確率を書く）．}
  \label{fig:hdrop}
\end{figure*}

\section{まとめと今後の展望}
本研究は ResTNet を動物将棋と 9×9 の将棋に適用し，Attention Map を中心に，IG で補強しながら，判断の説明に使えるかを確かめた．
完全解析された動物将棋では，T ブロックだけの 6T の 1 つのヘッド（第 5 T ブロックの第 4 ヘッド）が，選んだ例だけでなく無作為な局面でも，取れる相手のライオンや最善手の行き先に重みを向け，その上位のマスは IG の上位のマスとも偶然より多く重なった．
一方，R ブロックを含む構成では重みがほぼ均等であった．動物将棋の 9 構成の全ヘッドは最善手の行き先への重みがほぼ均等で，9×9 の将棋の R3(RRT) も重み × 81 の最大が均等の 1.03 倍であった．
9×9 の将棋では，図ごとに [0, 1] に正規化した Attention Map に濃淡は現れるが，実際に指された手の行き先や相手の玉のマスは濃くなく，手の確率を押し上げた駒（動かす銀，持ち駒の桂）は IG で読み取れた．
持ち駒は，動物将棋の 6T では有無で打つ先から重みを向けるマスが変わり，9×9 の将棋の R3(RRT) では Attention Map にほとんど表れなかった．示した局面では，どちらのゲームでも，打つ手の確率を最も押し上げたのは打つ持ち駒であった．

本研究には次の限界がある．
Attention Map が示すのは，起点の特徴を作るときにどのマスの情報を用いたかであり，モデルが手を選んだ理由そのものではない．
例として示した局面は説明しやすいものを選んでおり，無作為な局面で正解のマスと照らし合わせたのは，動物将棋の 6T と RRTRRT の，手番側のライオンを起点とした重みだけである．
9×9 の将棋は 1 局の 3 局面で，正解がないため，Attention Map と IG が将棋として正しいかは確かめられていない．

今後の展望として，次が挙げられる．
\begin{itemize}
  \item 人間の棋譜の多くの局面で，強いエンジンの最善手と Attention Map・IG を照らし合わせ，9×9 の将棋全体で言えることを調べる．
  \item R ブロックを含む構成で重みが均等になる理由を，学習の設定や構成を変えて調べる．
  \item IG の値が完全解析の正解と合うかを，動物将棋の全局面で確かめる．
\end{itemize}

\begin{thebibliography}{99}
\bibitem{alphazero} Silver, D., Hubert, T., Schrittwieser, J., Antonoglou, I., Lai, M., Guez, A., Lanctot, M., Sifre, L.,
  Kumaran, D., Graepel, T., Lillicrap, T., Simonyan, K.\ and Hassabis, D.: A General Reinforcement Learning Algorithm that
  Masters Chess, Shogi, and Go through Self-Play, \textit{Science}, Vol.~362, No.~6419, pp.~1140--1144 (2018).
\bibitem{restnet} Ju, Y.-R., Wu, T.-L., Shih, C.-C.\ and Wu, T.-R.: Bridging Local and
  Global Knowledge via Transformer in Board Games, \textit{Proc.\ IJCAI}, pp.~7446--7454 (2025).
\bibitem{gumbel} Danihelka, I., Guez, A., Schrittwieser, J.\ and Silver, D.: Policy
  Improvement by Planning with Gumbel, \textit{Proc.\ ICLR} (2022).
\bibitem{minizero} Wu, T.-R., Guei, H., Peng, P.-C., Huang, P.-W., Wei, T.-H., Shih, C.-C.\ and Tsai, Y.-J.:
  MiniZero: Comparative Analysis of AlphaZero and MuZero on Go, Othello, and Atari Games,
  \textit{IEEE Transactions on Games}, Vol.~17, No.~1, pp.~125--137 (2025).
\bibitem{ig} Sundararajan, M., Taly, A.\ and Yan, Q.: Axiomatic Attribution for Deep Networks,
  \textit{Proc.\ ICML}, PMLR, Vol.~70, pp.~3319--3328 (2017).
\bibitem{mandai2018} 万代悠作，金子知適：囲碁ニューラルネットワークの判断根拠の可視化，
  ゲームプログラミングワークショップ 2018 論文集，pp.~9--15 (2018).
\bibitem{ota2024} 大田皐介，中沢実：グラフ表現による将棋 AI の指し手の判断根拠の可視化手法，
  ゲームプログラミングワークショップ 2024 論文集，pp.~9--16 (2024).
\bibitem{tanaka2009} 田中哲朗：「どうぶつしょうぎ」の完全解析，
  情報処理学会研究報告ゲーム情報学（GI），Vol.~2009-GI-22，No.~3，pp.~1--8 (2009).
\bibitem{demura2023} 出村洋介，金子知適：将棋 81 万：強化学習のための多様性を持った将棋初期局面集，
  ゲームプログラミングワークショップ 2023 論文集，pp.~111--118 (2023).
\bibitem{aoba2020} 山下宏，保木邦仁，小林祐樹：AobaZero の紹介と観戦記，
  コンピュータ将棋協会誌，Vol.~31，pp.~38--44 (2020).
\bibitem{aoba2021} 山下宏，保木邦仁，小林祐樹：AobaZero の高速化と現在の状況，
  コンピュータ将棋協会誌，Vol.~32，pp.~72--86 (2021).
\end{thebibliography}

\end{document}
